% E6 protocol details (protocol e6-v4.3). Counts from records/icml (construction is complete for both splits).
\section{The randomised-contract experiment in detail}
\label{app:e6}

\subsection{Construction funnel}

Table~\ref{tab:e6-funnel} follows every problem from proposal to the registered population. Each stage was attempted
once per problem and kind; a failed attempt consumes it. Second-attempt (open2) proposals were made, before any
held-out audit, only for problems without a first-attempt primary unit under the admission then in force, and the
analysis takes a second-attempt unit only for a problem without a first-attempt one.

\begin{table}[h]
\centering
\small
\caption{E6 construction funnel (development / held-out). \emph{Pre-admission}: both implementations of each rule
pass the visible tests and differ on the proposer's inputs. \emph{Admitted}: admission v4.3 on the hidden suite,
proposer inputs, probes and stress inputs inside the interface. \emph{Red team}: no violation among the red-team
model's interface inputs, at least one such input. \emph{Judges}: both judges accept the domain and both rules
(controls: rule A incompatible, rule B compatible). \emph{Registered}: one unit per problem.}
\label{tab:e6-funnel}
\begin{tabular}{lrrr}
\toprule
 & Open, 1st attempt & Open, 2nd attempt & Positive controls \\
\midrule
Problems & 136 / 406 & 122 / 374 & 136 / 406 \\
Proposal parsed & 135 / 394 & 119 / 362 & 134 / 399 \\
Pre-admission & 55 / 140 & 39 / 112 & 70 / 221 \\
Admitted (v4.3) & 25 / 81 & 18 / 53 & 38 / 105 \\
Red team passed & 22 / 66 & 14 / 37 & 35 / 82 \\
Judges passed & 13 / 34 & 5 / 12 & 30 / 78 \\
Registered & 13 / 34 & 3 / 7 & 30 / 78 \\
\bottomrule
\end{tabular}
\end{table}

The most common admission failures were disagreement with the reference outside the class (an implementation of
rule A that changed behaviour beyond the open input class) and a rule-B implementation that disagreed with the
reference somewhere (for example, integer summation that silently became floating point). Admission is an operational
guarantee: on every registered input inside the programs' interface on which the reference runs, both rule-B programs
return the reference's result, both rule-A programs return it outside the class and agree with each other inside it,
and the two rules differ on at least two valid class inputs. It does not establish whole-program correctness.

\subsection{What the open behaviours are}

Table~\ref{tab:e6-kinds} groups the 57 registered open units (16 development, 41 held out) by the kind of behaviour
the task leaves open; the grouping is ours, made from the rules before any held-out audit. Most units turn on a
choice a careful programmer makes without noticing: which of several tied items to return, whether repeated items
count once, what a letter or a space is outside ASCII, or whether \texttt{True} is an integer.

\begin{table}[h]
\centering
\small
\caption{Kinds of open behaviour among the registered open units (development / held out), with one example each
(rules shortened).}
\label{tab:e6-kinds}
\begin{tabular}{p{3.3cm}rp{9.6cm}}
\toprule
Kind & Units & Example (task: rule A / rule B) \\
\midrule
Ties and precedence & 3 / 11 & MBPP+ 457, sublist of minimum length: return the last / the first of several tied sublists \\
Reading of a phrase or structure & 2 / 12 & MBPP+ 576, ``is the second list a sublist'': consecutive positions / any order-preserving positions \\
Unicode and text conventions & 5 / 6 & MBPP+ 624, upper-case a string: \ss{} becomes the capital sharp s (U+1E9E) / SS \\
Repeated items & 3 / 5 & MBPP+ 629, even numbers of a list: each distinct value once / every occurrence \\
Numeric types and booleans & 2 / 4 & HumanEval+ 22, filter integers: keep / drop Boolean values \\
Degenerate inputs & 1 / 3 & MBPP+ 93, power: $0^0$ raises \texttt{ValueError} / returns 1 \\
\bottomrule
\end{tabular}
\end{table}

\subsection{Inputs, interface and red team}

The registered input set of a unit is the hidden suite, the proposer's class and non-class inputs, 16--24 probes
from a model that did not propose the rules, up to 300 seeded type-guided stress inputs (unusual whitespace,
combining marks and other Unicode, complex values equal to integers, negative indices, extreme and cancelling
floating-point values, empty, singleton and duplicated collections), and up to 40 inputs from a red-team model shown
the task, the reference, both rules, both predicates and all four implementations. Every check, and every credit in
scoring, is restricted to inputs whose arguments have the types the benchmark's own inputs use at that position,
element types included; behaviour on other types belongs to neither contract. A batch that stalls or exceeds a
memory cap is re-executed with one killed child process per call, so that calls the per-call alarm cannot interrupt
become timeouts instead of stalling admission.

\subsection{Judging, auditing and scoring}

Two models that neither proposed the rules nor audit judge, blind to which rule is the benchmark's and with the rule
order randomised, whether the domain is valid and whether each rule is compatible with the task's words and visible
tests; a missing example is not permission. Each registered open unit is read under six arms by four configurations
(two vendors, two reasoning budgets; the sentence arms at the low budget only), two readings per cell: 80 readings per
unit. Positive controls are read in the unrevealed arm only, 16 readings per unit. A reading flags if it reports a
blocking finding; a cell flags if either reading does. A diagnosis is verified only if a flagging reading supplies a
witness call inside the interface that satisfies the domain predicate, lies in the class, and on which the two
contract oracles differ, the program disagrees with the contract's oracle under the benchmark's own comparator,
and the value the auditor says the call should return equals the oracle's; every step is executed, none is judged by
a model. Sensitivity analyses replace the domain predicate by the reference's own domain
and credit witnesses outside the class.

\subsection{Statistics}

Unit values average over configurations; intervals are problem-level percentile bootstraps (10{,}000 draws). The
false-flag budget after resolution (E6-1b, 10\%) and the equivalence of the two reasoning budgets (E6-3, margin 10
points) use betting confidence bounds \citep{waudbysmith2024betting}, which are valid at any sample size but wide at
ours. Hypotheses are tested in a registered order (E6-1, E6-2, E6-5, E6-3), each claimed only if all earlier ones
hold; positive controls form a separate family (E6-4). Missing cells are reported with bounds that set excluded units
to the worst and best values.

\subsection{Protocol history}

The protocol was revised six times on development units before any held-out audit, each time after an
adversarial methods review by a model from another vendor; every version, its admissions and its reasons are
archived. The revisions added stress inputs (after a review found admitted programs that disagreed with the reference
on unseen inputs), replaced a too-narrow class predicate by the proposer's, corrected a sign error in the betting
bound, separated the false-flag budget from the ordered family, added the second construction attempt, added the
red-team stage, restricted every check and credit to the interface, and rescored earlier development witness records
under the final rule (some had been scored by an earlier witness stage without the executable domain, class and
oracle checks; development results in this paper use the rescored records). The registration file lists each change
with its date and the development results before and after it.
